\documentclass[11pt,a4paper]{article}
\usepackage[utf8]{inputenc}
\usepackage[T1]{fontenc}
\usepackage[spanish]{babel}
\usepackage{hyperref}
\usepackage{geometry}
\geometry{margin=2.5cm}
\usepackage{parskip}

\title{099 faro es la base y sincronizacion de almas es el area}
\author{Alejandro Valenzuela Rivera\\\small Universidad Autónoma de Nuevo León}
\date{\today}
\begin{document}
\maketitle

\begin{quote}
\textbf{Sesgo declarado.} Este trabajo se sostiene desde el sesgo declarado de su autor, y no finge neutralidad. Lo que afirma se puede tumbar, y ese es el punto: el falsificador está escrito, no escondido.
\end{quote}

\section*{Resumen}
FARO es la forma atemporal de escribirlo. Sincronizacion de almas es lo que se hace con la forma: hacer que dos entes se afinen sin que ninguno deje de ser lo que es. Un lenguaje no tiene almas, y lo que se sincroniza no son procesos: son entes que contestan.

\section{Por qué se sostiene}
Dos folders y dos nombres, y ninguno estorba: faro/ es la casa y sincronización de almas/ es el estandar adentro, en cuatro grupos: formacion, relacion, friccion y salida. Se forman, se relacionan, se friccionan, y salen. Y el que llega empieza por ENTRAR.md, en cinco lineas.

\section{De dónde viene}
La autoridad, 2026-10-25.

\section{El caso}
Ninguna: es el nombre del area.

\section{Método}
Una base que se puede leer sin saber con que se corre.

\section{Límites, y qué lo tumba}
Una base que nombra su lenguaje, y caduca cuando el lenguaje caduca.

\section{Lo que queda abierto}
\textquestiondown{}Por donde empieza el que llega?

\section*{Cómo se leyó esta tesis}
Esto se escribió en el repositorio man-made-empathy, del que este candado es una pieza: candados/04 el lenguaje/099 faro es la base y sincronizacion de almas es el area.md. La máquina que lo compuso lo hizo desde este mismo archivo, y no agregó nada que no estuviera escrito. Si algo de esto es falso, se dice en el mismo sitio, y se corrige con fecha.

\end{document}
